\cotefile{75}
\scanpage{26}
c) Orientation d'un ensemble fini $I$\inserted{(cardinal $n$).} La donnée de $I$ équivaut à celle d'un torseur
\[
 P=\operatorname{Isom}(\deleted{\unsure{[1,n]_{\mathbb Z}|[1,n]_{\mathrm{ent}}}}I_n,I)=\unsure{\operatorname{numérot}|\operatorname{numérot.}}(I)
\]
$(I_n\mathrel{\unsure{=|:}}\{m\in\mathbb N^{*}\mid m\leq n\}=\{\deleted{\unsure{0}}1,\unsure{2},\ldots,n\})$ sous le groupe
\[
 \deleted{\unsure{\mathfrak S|S}}\mathfrak S_n=\mathfrak S_{I_n}=\operatorname{Aut}(I_n).
\]

On a un hom canonique ``signature''
\diagram{\begin{tikzcd}[ampersand replacement=\&]
\mathfrak S_n \arrow[r,"s=\mathrm{sgn}"] \& \{\pm1\} \&[-1em] \text{noyau }\mathfrak A_n
\end{tikzcd}}
et par définition l'\unsure{ens} des 2 orientations de $I$ est
\[
 \unsure{\omega|w}_I=\deleted{\unsure{\operatorname{Isom}(|\operatorname{Isom}}}P\times_{\mathfrak S_n}\{\pm1\}\simeq P/\mathfrak A_n .
\]

\uline{Compatibilités avec} a) \uline{et} b)

\uline{Avec} b) On a sur $E_{(I,k)}=k^{I}$ un produit scalaire canonique
\[
 (x,y)=\sum_{i\in I}x_i y_i
\]
pour lequel $(e_i)$ est une base orthonormale. Ainsi $\deleted{\unsure{E}}I$ \deleted{\unsure{donne}}\inserted{\unsure{donne}} naissance à un module quadratique $E(I,\unsure{x|k|X})$. Ceci posé, on \inserted{définit} une \unsure{bijection} canonique
\[
 \omega_{E(I,x)}\simeq\omega_I
\]
\margin{Si $i_1,\ldots,i_n$ est une \unsure{numérotation|numération} de $I$, $e_{i_1}\ \unsure{\cdots|\text{--}}\ e_{i_{\unsure{n}}}$ est une base \inserted{\unsure{ord.}} \unsure{orthon.} de $E$, \ldots}
par transitivité dans
\diagram{\begin{tikzcd}[ampersand replacement=\&,column sep=large]
 \mathfrak S_n \arrow[r,hook,"\mathrm{incl.}"] \arrow[rr,bend right=24,"\mathrm{sgn}" near start] \& \mathrm{O}(\unsure{n|N},k) \arrow[r,"\det"] \& \{\pm1\}
\end{tikzcd}}

En particulier, on trouve (\unsure{comp.} avec a).),
\[
 \omega_{E(I,\mathbb R)}\simeq\omega_I
\]
% Arrow from the preceding low E(I,R) toward the note along the foot of the sheet.
\margin{\tikz[baseline=-.35ex]{\draw[->] (0,.12) .. controls (.17,.08) and (.31,-.06) .. (.49,-.08);}\,\unsure{considéré} comme \unsure{variété} \unsure{via|vide} \unsure{ses pts|ses points|supra}}
% archive sheet number: 25
